\documentclass[aps,prb,onecolumn,nofootinbib,superscriptaddress]{revtex4-2}

\usepackage{amsmath,amssymb,amsthm,mathtools,bm}
\usepackage{booktabs}
\usepackage[colorlinks=true,linkcolor=blue,citecolor=blue,urlcolor=blue]{hyperref}
\usepackage{microtype}

\newcommand{\I}{\mathbf{1}}
\newcommand{\E}{\mathbb{E}}
\newcommand{\Tr}{\operatorname{Tr}}
\newcommand{\rank}{\operatorname{rank}}
\newcommand{\diag}{\operatorname{diag}}
\newcommand{\sgn}{\operatorname{sgn}}
\newcommand{\cA}{\mathcal{A}}
\newcommand{\cF}{\mathcal{F}}
\newcommand{\cH}{\mathcal{H}}
\newcommand{\cI}{\mathcal{I}}
\newcommand{\cL}{\mathcal{L}}
\newcommand{\cV}{\mathcal{V}}
\newcommand{\bxi}{\boldsymbol{\xi}}
\newcommand{\ket}[1]{\lvert #1\rangle}
\newcommand{\bra}[1]{\langle #1\rvert}
\newcommand{\ketbra}[2]{\lvert #1\rangle\!\langle #2\rvert}

\begin{document}

\title{Tangent-Cocycle Diagnostics for a Chiral Monitored Domain Wall}
\author{Asadullah Bhuiyan}
\date{\today}

\begin{abstract}
This note defines the numerical object used to test whether a topological
domain wall in the class-A adaptive fermion circuit carries a coherent chiral
channel.  The object is the differential of the \emph{normalized,
record-resolved} Gaussian state update, including both projective measurement
and conditional gain/loss reset.  It is not the singular-value spectrum of a
cumulative Kraus operator, a Choi-covariance transfer chart, or the
trajectory-averaged channel.  We give the exact fixed-record differential,
its stable QR restriction to two predetermined Hamiltonian edge modes, and
observables for survival, isotropy, wall retention, directed propagation, and
leakage of the encoded perturbation into the classical record.  We then state
a preregistered numerical campaign and a staged claim criterion.  Passing the
tangent-channel tests would motivate a separate coherent-information erasure
test; it would not by itself establish a quantum error-correcting code.  No
numerical result is asserted here.
\end{abstract}

\maketitle
\setcounter{tocdepth}{2}
\tableofcontents

\section{Question and object dictionary}
\label{sec:question}

The adaptive preparation protocol of Ref.~\cite{BhuiyanPanJian2026} is a
physical quantum instrument: the state is measured in a spatially structured,
overcomplete one-particle frame and a record-dependent gain or loss operation
corrects undesired occupations.  A domain wall between dynamical phases
therefore poses an operational question.  Does the complete
measurement--correction cycle transmit a two-dimensional space of
single-particle perturbations along the wall without learning which
perturbation was injected?

Ordinary square transfer matrices are poorly conditioned for this question.
Exact projectors and gain/loss operations create null directions; in a Choi
particle or hole chart these appear as zero or infinite rapidities.  Those
endpoints are meaningful data about a singular Gaussian operator, but they can
obscure the finite channel that acts on a state perturbation.  The differential
of the physical covariance update instead follows the sensitivity of the
normalized trajectory itself.  It is finite on every nonzero-probability
branch, although it may be rank deficient.  Table~\ref{tab:objects} records the
distinctions that will be maintained throughout.

\begin{table}[h]
\caption{\label{tab:objects}Related but inequivalent transfer objects.}
\begin{ruledtabular}
\begin{tabular}{p{0.23\linewidth}p{0.34\linewidth}p{0.35\linewidth}}
Object & Operational question & Present use \\
\hline
Choi/Kraus transfer data & What singular sectors does the realized operator
$K_{\bxi}$ possess? & Formal representation of projections and gain/loss; not
the numerical definition below. \\
Tangent cocycle & How does $G_0+\epsilon\,\delta G_0$ change along the same
record? & Primary two-mode transmission diagnostic. \\
Mean channel & How does $\sum_{\bxi}K_{\bxi}\rho K_{\bxi}^\dagger$ evolve
after discarding the record? & Control only; it does not determine quenched
trajectory sensitivity. \\
Flagged instrument & Can a record-dependent recovery preserve arbitrary
superpositions? & Eventual error-correction question, tested only after the
tangent gate. \\
\end{tabular}
\end{ruledtabular}
\end{table}

In particular, the quenched product of branch differentials cannot be replaced
by the differential of an outcome-averaged map.  Nor does a slowly contracting
tangent direction automatically define a logical degree of freedom: it may be
a classical metastable coordinate that is continuously revealed by the
counting record.

\section{Exact fixed-record branch differential}
\label{sec:differential}

\subsection{Measurement and reset}

Write the number-conserving covariance as
$G=2\langle c^\dagger c\rangle-\I$, and let
$P=\ketbra{\chi}{\chi}$ and $Q=\I-P$ for a normalized measured mode.
The outcome $s=+1$ denotes occupied and $s=-1$ empty.  Its conditional Born
probability and normalized covariance update are
\begin{align}
 p_s(G)&=\frac{1+s g}{2},
 &g&=\bra{\chi}G\ket{\chi},
 \label{eq:born}\\
 F_s(G)&=sP+Q(\I+sGP)^{-1}GQ .
 \label{eq:measurement-map}
\end{align}
For a \emph{fixed} outcome $s$, differentiation gives
\begin{align}
 D F_s(G)[\delta G]
 &=L_s(G)\,\delta G\,L_s(G)^\dagger,
 &L_s(G)&=Q(\I+sGP)^{-1}.
 \label{eq:measurement-differential}
\end{align}
One may multiply $L_s$ by an overall phase; in particular, the minus sign used
by a one-particle implementation has no effect on the covariance congruence.
Equation~\eqref{eq:measurement-differential}, rather than a finite difference
between two independently sampled outcomes, is the branch differential.  The
discrete draw $s$ is held fixed.

A conditional swap with an ancilla resets the same mode to covariance
$r=+1$ (filled) or $r=-1$ (empty):
\begin{align}
 R_r(G)&=QGQ+rP,
 &D R_r(G)[\delta G]&=Q\,\delta G\,Q .
 \label{eq:reset-differential}
\end{align}
Thus a corrected event has the column factor
\begin{equation}
 L_{s,r}^{\mathrm{event}}(G)=Q L_s(G).
 \label{eq:event-factor}
\end{equation}
For an ideal rank-one measurement and a reset of exactly the same normalized
mode, $QL_s=L_s$.  Omitting a separate $Q$ multiplication in this specific
branch therefore gives the same Jacobian in exact arithmetic; legacy GPU
spectra are not invalidated for this reason alone.  New code and tests
nevertheless represent the reset explicitly, both to audit chronology and to
avoid assuming this identity for generalized supports or imperfect mode
matching.

Equations~\eqref{eq:measurement-map}--\eqref{eq:event-factor} are evaluated
only when $p_s(G)>0$.  An exactly forbidden forced branch is invalid; replacing
its singular denominator by a numerical regulator and calling the result a
physical Lyapunov step changes the problem.  Production data must therefore
save the minimum Born denominator, null ranks, and the number of invalid forced
branches.

\subsection{Chronology and covariance perturbations}

Let $L_{t,j}$ be the factor for event $j$ of cycle $t$, including the reset
when it occurs, and define the chronological cycle factor
\begin{equation}
 \cL_t=L_{t,M_t}\cdots L_{t,2}L_{t,1}.
 \label{eq:cycle-factor}
\end{equation}
For a realized record $\bxi$, the restricted one-particle product and the full
covariance tangent are
\begin{align}
 \cA_T(\bxi)&=\cL_T\cdots\cL_2\cL_1,
 \label{eq:chronological-product}\\
 \delta G_T&=\cA_T\,\delta G_0\,\cA_T^\dagger .
 \label{eq:covariance-tangent}
\end{align}
This convention puts the factor acting on perturbation columns on the left.
It is equivalent, after adjointing the one-particle transfer convention, to
the usual Gaussian Schur-complement differential.  Its Lyapunov exponents are
quenched singular-value growth rates of a state derivative, not many-body
excitation gaps and not eigenvalues of an averaged Lindbladian.

\section{A QR-stabilized physical edge channel}
\label{sec:edge-channel}

\subsection{Predetermined edge frame}

The input frame is chosen before any monitored trajectory is generated.  Let
$H_{\mathrm{DW}}$ be the exact or spectrally flattened single-particle domain-
wall Hamiltonian used to define the protocol.  On a chosen wall, take two
orthonormal edge-band modes $f_0,f_1$ at adjacent momenta nearest the edge
crossing and set
\begin{equation}
 W_0=(f_0,f_1),\qquad W_0^\dagger W_0=\I_2.
 \label{eq:initial-frame}
\end{equation}
The localization and energy-window rules, including tie breaking, are fixed
once for every system size.  The two modes are not learned from the slowest
Lyapunov vectors.  This prevents the diagnostic from manufacturing a channel
by optimizing its input on the same disorder records used for evaluation.

Directly multiplying $\cA_TW_0$ underflows or overflows.  Starting from
$Q_0=W_0$ and $B_0=\I_2$, perform one thin QR factorization per full cycle,
\begin{align}
 Y_t&=\cL_t Q_{t-1}=Q_tR_t,
 &B_t&=R_tB_{t-1}.
 \label{eq:restricted-qr}
\end{align}
Then the identity
\begin{equation}
 \cA_tW_0=Q_tB_t
 \label{eq:factorization-identity}
\end{equation}
is exact up to roundoff.  Stabilize the $2\times2$ coordinate matrix by writing
$B_t=e^{\Gamma_t}\widehat B_t$, adding a scalar log norm to $\Gamma_t$ after
each multiplication.  Singular values are reconstructed as
$\log\sigma_i(B_t)=\Gamma_t+\log\sigma_i(\widehat B_t)$.  An exact zero
singular value is recorded as a null channel, not silently clipped into a
finite exponent.

\subsection{Survival, isotropy, and spatial support}

Let $\sigma_1\geq\sigma_2\geq0$ be the singular values of $B_T$.  The common
survival exponent and the scale-free anisotropy are
\begin{align}
 \overline\lambda_T
 &=\frac{\log\sigma_1+\log\sigma_2}{2T},
 \label{eq:survival}\\
 \delta_{\mathrm{iso}}(T)
 &=\frac{\sigma_1^2-\sigma_2^2}{\sigma_1^2+\sigma_2^2}
 =\tanh\!\left|\log\frac{\sigma_1}{\sigma_2}\right|.
 \label{eq:isotropy}
\end{align}
A neutral but strongly anisotropic channel preserves at most one direction;
it cannot carry an arbitrary two-component perturbation.  Both quantities are
invariant under a $U(2)$ change of the initial edge basis.

For a real-space projector $\Pi_w$ onto a fixed-width strip around the chosen
wall, define
\begin{align}
 \eta_w(t)&=\frac{\|\Pi_w Q_tB_t\|_F^2}{\|Q_tB_t\|_F^2},
 &\eta_{w\rightarrow\bar w}(t)
 &=\frac{\|\Pi_{\bar w}Q_tB_t\|_F^2}{\|Q_tB_t\|_F^2}.
 \label{eq:wall-retention}
\end{align}
The common scalar $e^{\Gamma_t}$ cancels.  The strip width is specified in
lattice units before the run and held fixed in the size extrapolation.  The
first ratio measures wall retention and the second leakage to the opposite
wall; their sum need not equal one because there is also intervening bulk.

\subsection{Directed propagation from interference}

Contraction alone does not establish chirality.  For adjacent momenta, inject
$u_+=(1,e^{i\phi_0})^T/\sqrt2$ and form the normalized propagated density
\begin{equation}
 \rho_t(x,y,\mu)=
 \frac{|[Q_t\widehat B_tu_+]_{x,y,\mu}|^2}
 {\|Q_t\widehat B_tu_+\|^2}.
 \label{eq:interference-density}
\end{equation}
Its first circular moment along the periodic direction is
\begin{equation}
 Z_t=\sum_{x,y,\mu}e^{2\pi i y/N_y}\rho_t(x,y,\mu),
 \qquad
 y_c(t)=\frac{N_y}{2\pi}\operatorname{unwrap}\arg Z_t .
 \label{eq:circular-center}
\end{equation}
The slope of $y_c(t)$ after burn-in defines a trajectory velocity.  The
analysis is repeated with several fixed $\phi_0$ values to exclude a phase-
choice artifact.  A chiral domain-wall interpretation requires nonzero mean
velocities with opposite signs on the two oppositely oriented walls, while
randomized and raster schedules quantify ordering bias.

\section{Does the counting record learn the edge coordinate?}
\label{sec:fisher}

Near a maximally mixed reference, a two-mode Hermitian perturbation can be
written
\begin{equation}
 \delta G_a(0)=W_0\sigma_aW_0^\dagger,
 \qquad a\in\{x,y,z\},
 \label{eq:logical-generators}
\end{equation}
where $\sigma_a$ are Pauli matrices.  Propagate each generator by the same
fixed record using Eq.~\eqref{eq:covariance-tangent}.  Immediately before
measurement event $j$, set
\begin{equation}
 d_{j,a}=\bra{\chi_j}\delta G_a(j^-)\ket{\chi_j},
 \qquad g_j=\bra{\chi_j}G(j^-)\ket{\chi_j}.
 \label{eq:score-data}
\end{equation}
The realized score and the conditional classical Fisher increment are then
\begin{align}
 h_{j,a}
 &=\partial_a\log p_{s_j}
 =\frac{s_jd_{j,a}}{1+s_jg_j},
 \label{eq:realized-score}\\
 J^{(j)}_{ab}
 &=\sum_{s=\pm1}p_s h_{j,a}(s)h_{j,b}(s)
 =\frac{d_{j,a}d_{j,b}}{1-g_j^2},
 \label{eq:fisher-increment}
\end{align}
for $|g_j|<1$.  Deterministic endpoints are evaluated by the corresponding
physical limit rather than by dividing by zero.  Because conditional score
increments form a martingale difference, the expected Fisher information of
the complete record is the expected sum of Eq.~\eqref{eq:fisher-increment}.
The principal statistic is its largest eigenvalue per observed cycle.

This quantity supplies the missing ``record privacy'' test.  A small
$\delta_{\mathrm{iso}}$ says that the normalized tangent map treats two input
directions similarly; it does not say that their Born weights are the same.
A neutral wall coordinate with finite Fisher-information density is therefore
a long-lived \emph{classical measured mode}, not a coherent memory.

For comparison, append the classical flag to a fixed-record instrument,
\begin{equation}
 \cI(\rho)=\sum_{\bxi}\ketbra{\bxi}{\bxi}\otimes
 K_{\bxi}\rho K_{\bxi}^\dagger.
 \label{eq:flagged-instrument}
\end{equation}
If recovery may depend on the accessible flag, exact Knill--Laflamme
correctability of a code projector $P_{\cV}$ requires, branch by branch,
\begin{equation}
 P_{\cV}K_{\bxi}^\dagger K_{\bxi}P_{\cV}
 =p_{\bxi}P_{\cV}.
 \label{eq:branch-kl}
\end{equation}
The record Fisher test probes the infinitesimal state-independence of
$p_{\bxi}$, while two-mode isotropy probes part of the normalized branch
action.  Neither test alone, nor their conjunction at finite size, proves
Eq.~\eqref{eq:branch-kl} or constructs a recovery map
\cite{KnillLaflamme1997}.

\section{Scattering interpretation and its limit}
\label{sec:scattering}

Separate common amplification from channel anisotropy by defining
\begin{equation}
 g_T=(\sigma_1\sigma_2)^{1/2},\qquad
 \widetilde\sigma_{1,2}=\sigma_{1,2}/g_T=e^{\pm x_T},
 \qquad
 x_T=\frac12\log\frac{\sigma_1}{\sigma_2}.
 \label{eq:rapidity}
\end{equation}
Then $\delta_{\mathrm{iso}}=\tanh(2x_T)$.  In a reciprocal flux-conserving
two-port transfer problem, the usual polar decomposition identifies
$\tau=\operatorname{sech}^2x_T$ as a transmission eigenvalue
\cite{ChalkerCoddington1988,Beenakker1997}.  This supplies a useful dictionary:
the QR cocycle measures normalized sensitivity transmitted from one temporal
boundary to the other, while $x_T$ is its imbalance rapidity.

The equality with a physical particle-transmission probability is
\emph{conditional}.  A monitored gain/loss tangent need not preserve a flux
metric or admit the reciprocal singular-value pairing of a unitary scattering
network.  Without an explicit current-conserving dilation, $\operatorname{sech}^2
x_T$ is a transmission \emph{proxy}, not an $S$-matrix probability.  Conversely,
the divergent directions of a singular Choi transfer chart should not be
inserted into Eq.~\eqref{eq:rapidity}: they encode exact occupied/empty caps of
the Kraus relation, whereas the present $\sigma_i$ describe the differential
of a normalized state trajectory.

\section{Numerical protocol and validation}
\label{sec:protocol}

The primary data are Born-sampled, perfect-correction trajectories of the
canonical adaptive circuit.  Postselection is retained only as a deterministic
implementation check.  The planned campaign is summarized in
Table~\ref{tab:campaign}; these are prospective parameters, not completed
results.  This note specifies intended behavior; implementation and completed-
run status must be established separately by source, tests, and output
metadata.

\begin{table}[h]
\caption{\label{tab:campaign}Prospective steady-ensemble campaign.}
\begin{ruledtabular}
\begin{tabular}{p{0.22\linewidth}p{0.70\linewidth}}
Quantity & Choice \\
\hline
Geometry & $N_x=20$, $N_y\in\{24,32,48,64,80\}$ \\
Frame truncation & $n_{\rm shell}=1$; repeat $N_y=32,64$ at $n_{\rm shell}=2$ \\
Time window & burn-in $2N_y$, observation $4N_y$ cycles \\
Domain-wall schedule & symmetric randomized ordering \\
Couplings & $\alpha_{\rm in}=1$, $\alpha_{\rm out}=30$, wall truncation on \\
Sampling & 128 trajectories for primary and uniform control; 64 otherwise \\
Primary feedback & Born-sampled perfect correction \\
Validation branch & forced postselection, never substituted for the primary ensemble \\
Controls & uniform topological, trivial hard wall with the same mask,
same-Chern artificial cut, and raster ordering \\
\end{tabular}
\end{ruledtabular}
\end{table}

Before data from this protocol are used for the staged claims of
Sec.~\ref{sec:gates}, the corresponding implementation must pass the following
tests.
\begin{enumerate}
 \item For a forced nonzero-probability branch, compare
 \begin{equation}
  \frac{F_{\bxi}(G+\epsilon\delta G)-F_{\bxi}(G-\epsilon\delta G)}{2\epsilon}
  \quad\text{with}\quad
  \cA_{\bxi}\delta G\cA_{\bxi}^\dagger
  \label{eq:finite-difference-test}
 \end{equation}
 over a convergence window in $\epsilon$.  Cover mixed and nearly pure $G$,
 occupied and empty outcomes, mismatch and no-mismatch branches, gain and
 loss, and local and dense update paths.
 \item Represent Eq.~\eqref{eq:reset-differential} explicitly in the CPU/GPU
 audit trail and test $QL_s=L_s$ for the current branch.  Under identical
 forced records, compare final covariance, event weights, QR factors, null
 counts, and restricted singular values.
 \item Check Eq.~\eqref{eq:factorization-identity} against an unstabilized
 small-system product, and verify independence under QR frequency and an
 arbitrary initial $U(2)$ rotation of $W_0$.
 \item Save denominator minima and censor invalid zero-probability branches
 rather than regularizing them into the physical sample.
\end{enumerate}

\section{Staged interpretation and claim gates}
\label{sec:gates}

The finite-size analysis uses trajectory bootstrap confidence intervals and
fits versus $N_y$ that retain raw null/censor counts.  The staged gate is:
\begin{enumerate}
 \item The magnitudes $|\overline\lambda_T|$, $\delta_{\mathrm{iso}}$, and the
 largest record-Fisher eigenvalue per cycle must each decrease with $N_y$ with
 a positive decay exponent at 95\% confidence.
 \item At $N_y=80$, require $\delta_{\mathrm{iso}}<0.1$ and a lower 95\%
 confidence bound on wall retention above $0.8$.
 \item On the two walls, the velocity confidence intervals must exclude zero
 and have opposite signs.
 \item No control may reproduce this entire conjunction.
\end{enumerate}
The numerical thresholds are conservative preregistration choices, not
universal constants.  Failure has a useful interpretation: depending on which
condition fails, the wall supports a localized slow tangent sector, a directed
but classically monitored perturbation, or a one-dimensional survivor, but not
the proposed coherent two-mode channel.

Only if all four conditions pass do we proceed to a Song-type erasure test for
a fixed-charge Gaussian candidate code
\begin{equation}
 \cV=\operatorname{span}\left\{
 f_0^\dagger\prod_{j=1}^{N-1}d_j^\dagger\ket{0},
 f_1^\dagger\prod_{j=1}^{N-1}d_j^\dagger\ket{0}
 \right\}.
 \label{eq:gaussian-code}
\end{equation}
Arbitrary superpositions in this two-dimensional space remain Slater
determinants because only one occupied orbital is varied.  For
$\overline\rho=P_{\cV}/2$ and a local erased region $A$, the coherent-
information loss is
\begin{equation}
 \Delta(A;\cV)=S_A(\overline\rho)+\log 2-S_{\bar A}(\overline\rho).
 \label{eq:song-loss}
\end{equation}
Its decay with distance and erased length gives a direct approximate-recovery
diagnostic in the spirit of chiral-edge and CFT codes
\cite{SongEtAl2026,SangHsiehZou2024}.  The mixture $\overline\rho$ is generally
non-Gaussian even though both codewords are Gaussian; its entropy cannot be
computed from the averaged covariance alone.  This gated calculation must be
validated against exact small-system Fock-space reductions.

This proposal is adjacent to dynamically generated and Floquet codes, where a
measurement schedule transports a genuine logical algebra and may enforce
anomalous boundary flow~\cite{HastingsHaah2021,AasenHaahLiMong2023}.  The
present adaptive circuit instead begins as dissipative state preparation.
Until Eq.~\eqref{eq:branch-kl}, coherent-information recovery, or an equivalent
operator-algebra test is demonstrated, the justified language is \emph{chiral
tangent channel} or \emph{boundary anomaly of active stabilization}, not
\emph{dynamical edge code}.

\section{Conclusion}

The proposed calculation replaces a singular transfer-core diagnostic by the
well-defined differential of the actual measurement--feedback trajectory.
Its QR restriction asks a deliberately narrow question: do two physically
chosen Hamiltonian edge modes survive, remain isotropic and wall-localized,
move chirally, and remain hidden from the Born record?  The combined answer is
more discriminating than a near-zero Lyapunov exponent.  It also creates a
clean decision boundary between three possible outcomes---a slow classical
wall coordinate, a coherent chiral tangent channel, and a genuine approximate
edge code---without presupposing that the strongest interpretation is true.

\bibliographystyle{apsrev4-2}
\bibliography{references}

\end{document}
